\documentclass[10.5pt]{config}
\usepackage{float}
\setlength{\parskip}{\baselineskip} % 设置段落间距为一个基本行距
\major{生物科学}
\name{蒋贤迪}
\partner{} % 签名栏
% \title{}
\stuid{3240105782}
\date{2025.10.22}
\lab{生物实验中心-312}
\course{生态学基础及实验}
\instructor{胡亮亮}
\grades{}
\expname{植物根形态性状的观测} % 实验名字
\exptype{观测} % 实验类型


\usepackage{amsmath} % 用于数学公式
\usepackage[table]{xcolor} % 用于设置单元格颜色
\usepackage{booktabs} % 用于更好的表格线条
\usepackage{array} % 提供 p{} 列类型和自定义列宽
\usepackage{booktabs} % 可选：提供更美观的表格线（此处未使用，但推荐）
\usepackage{makecell} % 支持多行文本和手动换行
\usepackage{caption} % 可选：用于标题
\usepackage{float}
\begin{document}


% \maketitlepage % 封面
% \maketoc % 目录
\makeheader % 首页头部


\section{实验目的和要求}
% --- 内容从 PDF 幻灯片 第8页 提取 ---
\textbf{实验目的：}

\begin{itemize}
\item 了解根的基本性状
\item 了解影响根性状的主要因素
\item 掌握获取及分析根性状的方法
\end{itemize}

\textbf{实验要求：}

\begin{itemize}
\item 描述植物种类、生长状况（可用照片反应）和生境（采样地点描述）
\item 统计描述不同植物各根性状指标并加以比较
\item 分析不同根性状指标间的关系
\item 利用 PCA 降维，判断植株样本之间是否形成某些聚类规律
\end{itemize}

\section{实验内容和原理}
% --- 内容从 PDF 幻灯片 第2-7页 提取 ---
\subsection{植物性状 (Plant Trait)}
植物从细胞到个体水平的任何可测量的\textbf{形态、生理或物候}特征，这些特征可能会影响其适应性或其环境。

\begin{itemize}
    \item 植物的\textbf{性状“套装” (a suite of traits)} 体现其\textbf{生存策略}
    \item 反映了在\textbf{资源获取 (Resource acquisition)} (如高光合速率、高比叶面积) 和\textbf{资源保守 (Resource conservatism)} (如叶片寿命长、组织密度高) 之间的\textbf{权衡 (Trade-off)策略}
\end{itemize}

\subsection{根性状 (Root Trait)}
\begin{itemize}
    \item \textbf{植物根系性状}是\textbf{功能性状}的重要组成部分，反映了植物如何投资生物量以获取地下资源。
    \item 不同植物的根千差万别
    \item 不同根性状之间存在许多关联
    \item 本实验主要讨论根性状中的\textbf{“根系构型（Root system architecture）”}
\end{itemize}

\textbf{双子叶和单子叶具有不同的根系构型}
        \begin{itemize}
            \item \textbf{主根系 (Tap root system)}: 通常在\textbf{双子叶植物 (Dicotyledon)}中发现。
            \item \textbf{须根系 (Fibrous root system)}: 通常在\textbf{单子叶植物 (Monocotyledon)}中发现。
        \end{itemize}


% --- 严格提取自 PDF 幻灯片 第5页 ---
\subsection{根性状的分类}
植物的根系是一个复杂的结构，为了系统地研究和比较不同植物的根系，生态学家建立了多种分类方案。这些方案从不同维度描述了根系的特征，主要可分为以下三类：

\subsubsection{基于发育与起源的分类}

这种分类根据根的发育来源和解剖学起源进行划分，是理解根系基本构型的基础。

\begin{itemize}
    \item \textbf{主根}: 由胚根直接发育而成，通常垂直向下生长，构成直根系的主轴
    \item \textbf{侧根}: 从主根或其他根内部中柱鞘细胞发育而成的分支
    \item \textbf{不定根}: 从茎、叶等非根器官发育而成的根
\end{itemize}

\subsubsection{基于根序的分类}

这是一种功能导向的分类法，将根系视为一个网络，从最细的根尖开始定义等级。

\subsubsection{基于直径的分类}

这是一种直观且功能意义明确的分类方法。

\subsection{影响根性状的环境因子}

\begin{itemize}
\item 土壤养分 Soil nutrients
\item 水资源分布 Water distribution
\item 邻体相互作用 Interaction with neighbors
\item 病害和毒害 Disease and toxicity
\item 共生微生物 Symbiotic microorganisms
\end{itemize}

根在响应环境的同时，也影响着环境

\subsection{移动性养分与非移动性养分}

前置知识：
\begin{itemize}
    \item 养分耗竭区是指养分能够被植物的根充分吸收消耗完的区域
\end{itemize}

\textbf{非移动性养分：}
\begin{itemize}
    \item 非移动性养分与土壤颗粒结合，根部很难获得。
    \item 养分耗竭区小。
    \item 更多根对于获取非移动性养分很有利,因为这扩大了吸收非移动性养分的范围
\end{itemize}

\textbf{移动性养分：}
\begin{itemize}
    \item 移动性养分在水中移动，因此根部很容易获得。
    \item 养分耗竭区大。
    \item 更多根对于获取移动性养分不算有利，因为这难以扩大吸收移动性养分的范围
\end{itemize}

\subsection{常见根性状指标}

\begin{minipage}[t]{0.6\textwidth}
    \subsubsection{几何信息}
    描述根系在空间中的整体分布范围：
    \begin{itemize}[leftmargin=*]
        \item \textbf{深度 (Depth)}: 根系垂直方向的最大距离
        \item \textbf{宽度 (Width)}: 根系水平方向的最大范围  
        \item \textbf{凸包面积 (Concex hull area)}: 包含整个根系的最小凸多边形面积
        \item \textbf{长度分布 (Length Profile)}: 根系在不同深度或宽度上的长度分布
    \end{itemize}
\end{minipage}%
\hfill
\begin{minipage}[t]{0.35\textwidth}
    \vspace{0.5em} % 手动调整垂直间距，使图片与文字对齐
    \centering
    \includegraphics[width=0.7\textwidth]{理论图/geometrical-info.png}
    \captionof{figure}{几何信息示意图}
\end{minipage}

\vspace{1cm}

% 第二组：形态信息
\begin{minipage}[t]{0.35\textwidth}
    \subsubsection{形态信息}
    描述单个根段的物理特性：
    \begin{itemize}[leftmargin=*]
        \item \textbf{直径 (Diameter)}: 根段的粗细程度
        \item \textbf{长度 (Length)}: 单个根段的长度
        \item \textbf{方向 (Orientation)}: 根段生长的方向角度
    \end{itemize}
\end{minipage}%
\hfill
\begin{minipage}[t]{0.35\textwidth}
    \vspace{0.5em}
    \centering
    \includegraphics[width=0.7\textwidth]{理论图/morphological-info.png}
    \captionof{figure}{形态信息示意图}
\end{minipage}

\vspace{1cm}

% 第三组：拓扑信息
\begin{minipage}[t]{0.6\textwidth}
    \subsubsection{拓扑信息}
    描述根系的分支结构和连接方式：
    \begin{itemize}[leftmargin=*]
        \item \textbf{侧根数量 (Lateral number)}: 父根上的分支数量
        \item \textbf{侧根间距 (Inter-lateral distance)}: 相邻分支之间的距离
        \item \textbf{分支角度 (Insertion angle)}: 侧根从父根长出的角度
        \item \textbf{未分支顶梢区长度 (Length of Unbranched Apical Zone, LAUZ)}: 根顶端未发生分支的长度
    \end{itemize}
\end{minipage}%
\hfill
\begin{minipage}[t]{0.35\textwidth}
    \vspace{0.5em}
    \centering
    \includegraphics[width=0.7\textwidth]{理论图/topological-info.png}
    \captionof{figure}{拓扑信息示意图}
\end{minipage}

此外，还有比较重要的拓补指标：

\begin{itemize}
        \item \textbf{海拔 (Altitude, \(a\))}: 从根基到最远根尖的连接数（最长路径长度）
        \item \textbf{量值 (Magnitude, \(μ\))}: 根系中所有外部连接（根尖）的总和
        \item \textbf{拓扑指数 (Topological index, TI)}: 计算为 \(\log a / \log μ\)，用于衡量根系连接效率
\end{itemize}

其中，对于海拔和量值，有以下补充：

\begin{minipage}[t]{0.6\textwidth}
    \textbf{根段海拔 (Altitude of Segment)}
    描述根段在根系网络中的层级位置：
    \begin{itemize}[leftmargin=*]
        \item \textbf{定义}: 从根系基部到该根段所需经过的连接数量
        \item \textbf{意义}: 衡量根段在分支结构中的"深度"，数值越大位置越末端
        \item \textbf{示例}: 如图示，从基部(海拔1)到最远端根段(海拔4)
    \end{itemize}
\end{minipage}%
\hfill
\begin{minipage}[t]{0.35\textwidth}
    \vspace{0em}
    \centering
    \includegraphics[width=0.4\textwidth]{理论图/altitude-diagram.png}
    \captionof{figure}{根段海拔示意图}
\end{minipage}

\vspace{1cm}

% 第二段：Magnitude of segment
\begin{minipage}[t]{0.6\textwidth}
    \textbf{根段量值 (Magnitude of Segment)}
    描述根段所支撑的下游分支规模：
    \begin{itemize}[leftmargin=*]
        \item \textbf{定义}: 从根尖开始向上游累积计算的数值
        \item \textbf{规则}: 根尖段量值为1，父根段量值为子根段量值之和
        \item \textbf{意义}: 根基部根段的总量值等于整个根系的根尖总数
        \item \textbf{示例}: 如图示，1+1=2, 2+2=4, 最终总量值=6
    \end{itemize}
\end{minipage}%
\hfill
\begin{minipage}[t]{0.35\textwidth}
    \vspace{0.5em}
    \centering
    \includegraphics[width=0.4\textwidth]{理论图/magnitude-diagram.png}
    \captionof{figure}{根段量值计算示意图}
\end{minipage}

\subsection{主成分分析 (Principal Component Analysis, PCA)}
\begin{minipage}[t]{0.6\textwidth}
    \textbf{PCA原理与应用}
    主成分分析是一种常用的多元统计方法：
    \begin{itemize}[leftmargin=*]
        \item \textbf{核心思想}: 将样本的$n$维特征数据映射到$l$维空间（$n>>l$），去除数据之间的冗余性。
        \item \textbf{算法动机}：保证样本投影后方差最大
        \item \textbf{降维目的}: 在保留大部分原始信息的同时，简化数据结构
        \item \textbf{根系分析}: 在根系研究中，PCA可用于识别影响根系构型的关键性状组合
        \item \textbf{结果解读}: 
            \begin{itemize}
                \item \textbf{PC1}: 通常代表影响最大的性状组合，解释数据中最多的变异
                \item \textbf{PC2}: 代表次重要的性状组合，与PC1正交且不相关
                \item \textbf{载荷}: 各原始变量对主成分的贡献程度
                \item \textbf{得分}: 样本在主成分空间中的位置
            \end{itemize}
    \end{itemize}
\end{minipage}%
\hfill
\begin{minipage}[t]{0.35\textwidth}
    \vspace{0.5em}
    \centering
    \includegraphics[width=0.95\textwidth]{PCA图/pca-diagram.png}
    \captionof{figure}{PCA分析结果示意图\\展示样本在主成分空间中的分布}
\end{minipage}

\section{实验材料与设备}
% --- 内容从 PDF 幻灯片 第9页 提取 ---

\begin{itemize}
\item 铲子
\item 扫描仪
\item 塑封袋
\item 根盘
\item RhizoVision Explorer 软件
    \item 剪刀和镊子
    \item 烘箱 (设置 $60^\circ\text{C}$)
\end{itemize}


\section{操作方法和实验步骤}
% --- 内容从 PDF 幻灯片 第10, 12, 14, 16页 提取 ---

\begin{enumerate}
    \item \textbf{植物样品采集}
    \begin{itemize}
        \item 在同一生境中，组内每位同学各自选定1个植物物种（禾本科、菊科等）
        \item 采集3个个体
        \item 完整挖出根部，保持细根完整
        \item 单独放入样本袋（包括地上部分和地下部分）
    \end{itemize}
    
    \item \textbf{根样清洗}
    \begin{itemize}
        \item 用自来水冲洗根样
        \item 完全去除表面杂物
    \end{itemize}
    
    \item \textbf{根系摆放}
    \begin{itemize}
        \item 根盘中加蒸馏水
        \item 仔细摆放根系，减少交叉
        \item 根系舒展，不交错重叠
        \item 保证托盘内无泥土颗粒等杂质
        \item 根和托盘四周留出1cm空隙
    \end{itemize}
    
    \item \textbf{根系扫描}
    \begin{itemize}
        \item 擦干根盘外部
        \item 扫描获取600 dpi图像
        \begin{itemize}
            \item 模式：照片模式
            \item 文稿来源：透扫适配器部件
            \item 文稿类型：彩色正片
            \item 图像类型：8位灰度
            \item 扫描质量：标准
            \item 图像格式：JPEG
        \end{itemize}
    \end{itemize}
    
    \item \textbf{数据分析}
    \begin{itemize}
        \item 使用RhizoVision Explorer软件
        \item 分析图像获取性状指标
    \end{itemize}
    
    \item \textbf{样品处理}
    \begin{itemize}
        \item 根样放入信封，写上上课时间、组号和植物名称
        \item 地上部分和地下部分放在同一信封中
        \item 60℃烘干48小时
    \end{itemize}
\end{enumerate}

\section{实验结果与分析}

\subsection{植物生长状况}

本次实验中，我们组采集了不少植物，我选择了其中的六种，并且为了客观获取并统计数据，我选择在“中国植物志”和“植物智”网站查询相关信息，这样能为后文的数据分析奠定基础。经过查询网站，得到如表\ref{tab:plant_info}所示的数据

\begin{table}[htbp]
\centering
\caption{植物生长状况}
\label{tab:plant_info}
\renewcommand{\arraystretch}{1.2} % 增加行高 (默认为1.0)
    \setlength{\tabcolsep}{8pt}      % 增加单元格左右内边距 (默认为6pt)
\begin{tabular}{|p{1.6cm}|p{3.5cm}|p{1.6cm}|p{1.6cm}|p{2.2cm}|p{5cm}|}
\hline
植物名称 & 拉丁学名 & 属 & 科 & 生活型 & 生长状态 \\
\hline
香附子 & \textit{Cyperus rotundus L.} & 莎草属 & 莎草科 & 多年生草本 & 具细长匍匐根状茎和黑褐色块茎；秆高15-95厘米，锐三棱形；叶基生，短于秆 \\
\hline
百日菊 & \textit{Zinnia elegans Jacq.} & 百日菊属 & 菊科 & 一年生草本 & 茎直立，高30-100厘米，被糙毛或长硬毛；叶对生，宽卵形或长圆状椭圆形，长5-10厘米\\
\hline
合欢 & \textit{Albizia julibrissin Durazz.} & 合欢属 & 豆科 & 落叶乔木 & 树冠开展，高达16米；二回羽状复叶，小叶昼开夜合；头状花序排成伞房状，花丝细长，粉红色\\
\hline
苎麻 & \textit{Boehmeria nivea (L.) Gaudich.} & 苎麻属 & 荨麻科 & 亚灌木或灌木 & 茎直立，高0.5-1.5米，多分枝，密被长硬毛；叶互生，长6-15厘米，叶背密被白色绵毛\\
\hline
麦冬 & \textit{Ophiopogon japonicus (L. f.) Ker-Gawl.} & 沿阶草属 & 天门冬科 & 多年生草本 & 根较粗，中部或近末端常膨大成椭圆形或纺锤形小块根；叶基生成丛，禾叶状，长10-50厘米\\
\hline
香堇菜 & \textit{Viola odorata L.} & 堇菜属 & 堇菜科 & 多年生草本 & 具匍匐茎；根状茎较粗，淡褐色；叶基生，圆形或宽卵状心形，长1.5-4.5厘米；花深紫色，有香气 \\
\hline
\end{tabular}
\end{table}

\newpage

其中，采集这些植物的生境具体信息和其他相关信息如图\ref{fig:habitat_images}和表\ref{tab:collected_plants}所示：

\begin{figure}[htbp]
    \centering
    
    % --- 第一张图片 ---
    \begin{minipage}{0.32\textwidth}
        \centering
        \includegraphics[width=\linewidth]{生境图/图1.jpg}
    \end{minipage}
    \hfill % 自动分配水平间距
    % --- 第二张图片 ---
    \begin{minipage}{0.32\textwidth}
        \centering
        \includegraphics[width=\linewidth]{生境图/图2.jpg}
    \end{minipage}
    \hfill % 自动分配水平间距
    % --- 第三张图片 ---
    \begin{minipage}{0.32\textwidth}
        \centering
        \includegraphics[width=\linewidth]{生境图/图3.jpg}
    \end{minipage}
    \caption{植物采样生境照片}
    \label{fig:habitat_images}
\end{figure}

\begin{table}[htbp]
\centering
\caption{采集植物的种类+生境}
\label{tab:collected_plants}
\begin{tabular}{lcccc}
\toprule
\textbf{植物名称} & \textbf{纲} & \textbf{目} & \textbf{科} & \textbf{生境} \\
\midrule
香附子 & 木兰纲 & 禾本目 & 莎草科 & 生物实验中学西南侧草地 \\
百日菊 & 木兰纲 & 菊目 & 菊科 & 生物实验中学南侧树林 \\
合欢 & 木兰纲 & 豆目 & 豆科 & 生物实验中学南侧树林 \\
苎麻 & 木兰纲 & 蔷薇目 & 荨麻科 & 生物实验中学南侧树林 \\
麦冬 & 木兰纲 & 天门冬目 & 天门冬科 & 生物实验中学南侧树林 \\
香堇菜 & 木兰纲 & 金虎尾目 & 堇菜科 & 生物实验中学南侧树林 \\
\bottomrule
\end{tabular}
\end{table}

\subsection{根系扫描处理}

\begin{itemize}
    \item 由于根系较为破碎，故本次的"Analysis mode"选择"broken roots"模式来分析数据。
    \item 由于拍摄获得的图像分辨率为600dpi，故在根系分析前勾选"Convet pixles to physical units"，并选择600dpi，相当于每英寸打印出600个点，经软件自动转换可得到正确的数据。
    \item 由于存在些许杂质，故勾选"Filter non-root objects"，通过调整Maximun size"，观察到当其值为3时效果最好。既可以剔除杂质，又不会剔除部分细根。
\end{itemize}

经Rhizo Vision Explorer软件分析根系图后，得到了以下根系扫描图：

\begin{figure}[htbp]
    \centering
    
    % --- 第一行图片 ---
    \begin{minipage}{0.3\textwidth}
        \centering
        \includegraphics[width=\linewidth]{植物根系图/香附子.png}
        \centerline{香附子}
    \end{minipage}
    \hfill % 自动分配水平间距
    \begin{minipage}{0.3\textwidth}
        \centering
        \includegraphics[width=\linewidth]{植物根系图/百日菊.png}
        \centerline{百日菊}
    \end{minipage}
    \hfill % 自动分配水平间距
    \begin{minipage}{0.3\textwidth}
        \centering
        \includegraphics[width=\linewidth]{植物根系图/合欢.png}
        \centerline{合欢}
    \end{minipage}

    \vspace{0.5cm} % 设置两行之间的垂直间距

    % --- 第二行图片 ---
    \begin{minipage}{0.3\textwidth}
        \centering
        \includegraphics[width=\linewidth]{植物根系图/苎麻.png}
        \centerline{苎麻}
    \end{minipage}
    \hfill % 自动分配水平间距
    \begin{minipage}{0.3\textwidth}
        \centering
        \includegraphics[width=\linewidth]{植物根系图/麦冬.png}
        \centerline{麦冬}
    \end{minipage}
    \hfill % 自动分配水平间距
    \begin{minipage}{0.3\textwidth}
        \centering
        \includegraphics[width=\linewidth]{植物根系图/香堇菜.png}
        \centerline{香堇菜}
    \end{minipage}

    \caption{六种代表性植物的根系扫描图}
    \label{fig:root_scans}
\end{figure}

其中，每张照片从左到右依次从1到3，每张图片的每张根系图下方都有标号。颜色可以代表根的直径，从红色到蓝色再到绿色，直径依次递增
\begin{itemize}
    \item 当根直径<2.00mm时，绘制成红色
    \item 当2.00mm<根直径<5.00mm时，绘制成蓝色
    \item 当根直径>5.00mm时，绘制成绿色
\end{itemize}

\newpage

\hspace{2em}观察图像，可发现：
\begin{enumerate}
    \item 香附子根系具明显块茎。样本1是密集的块茎和须根团，样本2和3则展示了连接块茎的细长根状茎。这种克隆生长结构使其根系形态与依赖主根的植物完全不同。香附子-3的根直观感觉少于1号和2号，可能是挖断根导致的。
    \item 百日菊的根都被挖断了，2、3号很明显，1号我判定挖断是因为在挖掘时真的挖断了，带来了一定误差。根系视觉上很特殊，样本均呈现出直径大、长度长的主根，形态弯曲。与其粗壮主根相比，侧根和须根数量显得稀少。
    \item 合欢的样本直观反映了木本采样困难。样本1非常破碎，仅为片段。样本2和3相对完整，展示了清晰的主根分支。破碎的样本为后续数据偏低埋下了伏笔
    \item 苎麻-2的根系的破碎程度各不相同，其中2号最为破碎，3号似乎完整但看上去并没有1号多，很有可能是挖断了根但未找到被挖断的根。其根系形态多变。其根系兼具粗壮的木质化特征和纤细须根，直观上体现了亚灌木介于草本和木本间的过渡形态。
    \item 麦冬根系呈放射状，具典型块根，根中段明显膨大。样本1和2构型完整，形态相似。而3号根明显不完整，这说明在挖掘时挖断了许多根。这种膨大结构在视觉上与其他草本差异显著，反映了其不同的养分储存策略。
    \item 香堇菜的根系具有膨大块根，用于水分和养分储存，且样本1和2的侧根稀疏且不完整，可能在挖掘时有较多断裂，而样本3相对更加完整
\end{enumerate}

\newpage

将所得数据整理后，得到如表\ref{tab:root_parameters_vertical}所示的数据

\begin{table}[htbp]
\centering
\caption{各项根性状数据}
\label{tab:root_parameters_vertical}
\renewcommand{\arraystretch}{1.5} % 增加行高
\setlength{\tabcolsep}{6pt}        % 调整单元格内边距
\begin{tabular}{lcccccc}
\toprule
植物名称 & 香附子 & 百日菊 & 合欢 & 苎麻 & 麦冬 & 香堇菜 \\
科属 & 莎草科 & 菊科 & 豆科 & 荨麻科 & 天门冬科 & 堇菜科\\
\midrule
平均根尖数量 & 2273.33 & 3074.33 & 1002.00 & 2317.67 & 3226.00 & 955.67 \\
平均总根长度 (mm) & 3080.41 & 4858.96 & 1536.50 & 2685.35 & 4335.56 & 1664.91 \\
平均网络面积 (mm²) & 1069.82 & 2698.81 & 474.92 & 702.55 & 2385.02 & 585.01 \\
平均直径 (mm) & 0.91 & 1.70 & 0.52 & 0.56 & 1.40 & 0.68 \\
平均周长 (mm) & 3275.45 & 3825.93 & 2329.23 & 3719.83 & 3177.92 & 2057.62 \\
平均体积 (mm³) & 7616.38 & 28740.99 & 1157.50 & 1162.37 & 11194.81 & 2166.46 \\
平均表面积 (mm²) & 9636.81 & 28585.73 & 2574.48 & 3951.67 & 18779.79 & 3767.29 \\
平均每毫米分支频率 &1.23 &0.82 &0.94& 1.15&0.79&1.17\\
分支点数量 & 3764.67 &4153.67 & 1400.67 & 3347.33 & 3564 & 1946.33\\

\bottomrule

\end{tabular}

\end{table}

在本次实验中，我选择了6种植物，其中包括四种草本植物，一种木本植物——合欢，和一种亚灌木——苎麻，其中亚灌木可以认为是草本向木本的过渡类型。

观察数据，可以发现：
\begin{enumerate}
    \item 在草本植物中，除去香堇菜，其他三种的多数指标值有一定差距，但少数指标相差并不大，如“平均根尖数量”和“平均总根长度”等。而香堇菜在各项指标中，都低于或远低于其他草本植物。我认为可能的原因有：
    \begin{itemize}
        \item 每个植物所处的生长状态并不同，如百日菊在各项指标中都处高位，是因为其已经比其他草本植物更加成熟，根系更加发达。
        \item 虽然都属于草本植物，但根与根之间亦有差距。从图中可直观感受到百日菊的根直径很大，而分析结果也体现了这点，其他指标亦是如此。
        \item 在挖掘香堇菜时，挖断了不少根数，导致其诸多指标都远低于同属于草本植物的其他植物，因此该实验中香堇菜的数值对整个草本植物的数值影响很大。
    \end{itemize}
    \item 属于木本植物的合欢，在所有指标中，都远低于草本植物（香堇菜除外）。造成这样的原因，我认为可能是因为：
    \begin{itemize}
        \item 我们所找到的合欢样本还处于幼年期，其根系不发达。
        \item 在取样品时，由于木本植物根系的特殊性，我们在挖掘时弄断的根数明显多于其他植物。如合欢-1，和其他样本相比，算是最不完整的根了。实际上合欢-3的很多项指标都是合欢-1的数倍，如根数、总根长、最大宽度.
    \end{itemize}
    \item 在诸多指标中，属于亚灌木的苎麻的数值都介于草本（除香堇菜）和木本之间，这也体现出亚灌木可以认为是草本与木本的过渡形态。
\end{enumerate}

\subsection{数据分析}

\subsubsection{\textbf{主成分分析（PCA）聚类}}



将表\ref{tab:root_parameters_vertical}数据进行标准化后，再进行PCA聚类，得到该示意图：

\begin{figure}[htbp]
    \centering
    \includegraphics[width=0.5\linewidth]{PCA图/PCA.png}
    \caption{PCA聚类示意图}
    \label{fig:PCA}
\end{figure}

观察该图，得知：
\begin{enumerate}
    \item 第一主成分$PC_1$的解释度为$68.53\%$，，第二主成分$PC_2$的解释度为$22.05\%$，二维总解释度已经到达了$90.58\%$，已经很好能够解释数据分布了。
    \item 其中部分样本分布较为密集，如紫色的香附子、橙色的合欢、棕色的香堇菜；而其余样本，如绿色的苎麻、红色的麦冬和蓝色的百日菊分布较为分散。
    \item 对于麦冬和百日菊，有两个样本间隔较短。第三个样本间隔较远，且偏远其他同类样本的\textbf{方向}和\textbf{距离}甚至都很相近。我个人认为：这可能是因为该两个样本的根系破损程度比较严重，因其根系的完整度和其他同类样本显著不同，故偏离方向和偏离距离都很相近。经过附录的得分表格可知，该异常的植株编号为百日菊-1，麦冬-3，且这两个样本的部分指标都明显低于其余两个样本，尤其是根数和根长，能够一定程度验证我的猜想。
    \item 对于苎麻，PCA聚类图呈现了很奇怪的分布特征，即三个点相隔较远且\textbf{几乎在同一条直线上}。对此，我认为：如果三个点距离远且无规律，那是采样问题；但如果在此基础上甚至可以回归到一条线上，那说明我们组所采的三个样本，很有可能处于不同的发育阶段，因此呈现出了很强的线性关系。并且，意识到这点后，我观察到其余多数样本也几乎落在同一直线上，也呈现了一定的线性关系，只有香堇菜为例外。当然，以上观点纯属个人想法，具体是否如此得等下次实验解封信封，并比对是否属于不同发育阶段后才能判断论点是否正确。
    \item 对于$PC_1$，它贡献了主要的解释度（$68.53\%$）。观察载荷箭头，绝大多数测量指标（如总根长、根尖数量、网络面积、平均直径、体积、表面积等）都与$PC_1$呈强正相关。这表明$PC_1$主要代表了根系的\textbf{“整体规模”}。$PC_1$值越大（越往右），代表根系越庞大、越粗壮（如百日菊和麦冬）；值越小（越往左），代表根系越小或越破碎（如合欢）。
    \item 对于$PC_2$，它贡献了$22.05\%$的解释度。其主要的正向载荷是“每毫米分支频率”。这表明$PC_2$主要反映了根系的\textbf{“分支策略”}。$PC_2$值越高（越往上），代表根系在单位长度上的分支越密集（如苎麻、香附子）；$PC_2$值越低，则分支频率越低。
    \item 针对“是否可以通过筛选优化分析结果”这一问题，答案是肯定的，并且可以从两个主要方面入手：\textbf{筛选样本}与\textbf{筛选性状}。
        \begin{itemize}\item \textbf{通过筛选样本进行优化：}如前文第2、3点的分析，当前PCA结果（尤其是$PC_1$）在很大程度上受到了\textbf{样本完整度}的干扰。这种“大小效应”掩盖了不同物种间真实的形态差异。因此，剔除那些明显破碎或不完整的样本（如麦冬-3和百日菊-1），仅对高质量的样本进行重新分析，是一种直接有效的优化方法，能更好地揭示相似发育阶段下不同物种的真实聚类规律。    
        \item \textbf{通过筛选性状进行优化：}另一种更深入的优化策略是筛选分析所用的性状。既然已知$PC_1$主要代表根系的“整体规模”，其贡献主要来自于总根长、网络面积、体积等强相关的“大小”指标。我们可以尝试进行一次新的PCA分析，在分析中有意排除有关“大小”的指标，采用有关分支模式、空间探索策略等内在形态特征上的性状指标。可惜的是，那些指标需要分析完整根时才能获得，故本实验难以优化PCA。
\end{itemize}
\end{enumerate}

\subsubsection{\textbf{层次聚类分析（Hierarchical cluster analysis）}}

通过层次聚类分析，选定所有性状，并指定簇数为6，可得到如图\ref{fig:all}结果：

\begin{figure}[htbp]
    \centering % 将整个大图居中

    % --- 这是左侧的子图 (图3) ---
    \begin{subfigure}[b]{0.43\textwidth}
        \centering
        % 将 "图片3的文件名.png" 替换为您的真实文件名
        \includegraphics[width=\textwidth]{树状/树状聚类.png}
        \caption{选择所有数据所得图}
        \label{fig:all} % 为子图设置一个独一无二的标签
    \end{subfigure}
    \hfill % 这个命令会自动添加水平方向的弹性空白，将两张图推开
    % --- 这是右侧的子图 (图4) ---
    \begin{subfigure}[b]{0.43\textwidth}
        \centering
        % 将 "图片4的文件名.png" 替换为您的真实文件名
        \includegraphics[width=\textwidth]{树状/without.png}
        \caption{剔除部分数据所得图}
        \label{fig:not all} % 为子图设置一个独一无二的标签
    \end{subfigure}

    \caption{层次聚类分析图} % 整个大图的总标题
    \label{fig:HCA} % 整个大图的总标签
\end{figure}

明显可见，该结果并不理想。我认为主要原因还是样本完整度的问题。例如：
\begin{itemize}
    \item 百日菊-1和百日菊-2、-3相隔甚远，很大原因是因为百日菊-1的根系并不完整。虽然百日菊-2、-3有断裂情况，但完整程度还是显著优于百日菊-1。
    \item 麦冬-3和麦冬-1、-2也能很好体现。从图\ref{fig:root_scans}可看出，麦冬-3的完整度显著比另外两个样本差，因此麦冬-1、-2能归为一类，而麦冬-3则不行
\end{itemize}

而中间有很多植物被归为一类，我推测主要是“大小”中的“体积”与“表面积”因素影响较大，因为：
\begin{enumerate}
    \item 当我依次添加性状指标时，我发现当不勾选“体积”和“表面积”时，没有出现这种现象.如图\ref{fig:not all}。
    \item 而当我添加任意其一后，均出现了一群植株被归为一类的情况，和\ref{fig:all}很相似
    \item 虽然分类效果并不理想，但至少没有出现那诡异的一大片粉色林。因此我认为粉色林中植株的根系，在“体积”和“表面积”方面呈现明显的相似性。
    \item 并且，我个人也难以对层次聚类进行调整。该网站不像PCA的网站，会展现出关键步骤，因此我不清楚关键这些关键性状是否经过标准化，且该网站也不允许我设定切割高度等其他设置，因而有了这张效果极不理想的图。
\end{enumerate}

\section{思考题}
\textbf{采集样本形成不同根系特征的原因是什么}

在演化进程中，不同植株选择不同的根系特征，是其祖先在不同环境中经受自然选择，从而形成独特生存策略的直接体现。例如，百日菊作为一年生草本，在演化中形成了粗壮的主根和稀疏的侧根，这是一种在养分水分移动性较强的环境中优先快速深扎以获取资源的策略。与之相反，香附子演化出块茎与须根的克隆生长结构，使其能在竞争激烈的草地生境中快速扩展并共享资源。合欢作为木本植物，其根系投资更倾向于构建深层稳定的锚固结构，以支持其高大的地上部分并实现长期生存，这与草本植物快速吸收的策略形成权衡，从图\ref{fig:root_scans}中就可看出合欢与其他草本植物有很大的区分度。而麦冬根中部的膨大块根，则是演化出的高效储藏器官，用于应对季节性的干旱或养分短缺。即便是作为草本向木本过渡的苎麻，其根系兼具较高的分支频率和初步的木质化，也恰恰反映了演化过程中的一种中间适应状态，例如\ref{fig:root_scans}中的苎麻-1与木本的合欢相似，而苎麻-3又和其他草本相似。
虽然本次实验有一些的样本被挖了，但幸运的是我们组有草本、木本和介于两者之间的亚灌木，形成了一定的根系多样性。这些根系是植物在亿万年演化史上，为适应各自特定的水分、养分、竞争压力等环境条件，而在“资源快速获取”与“长期生存保障”之间做出的进化权衡结果，其独特的根系特征正是写入基因的适应性解决方案。

\section{讨论、心得}

该实验给我留下很深印象，我将从两方面讲述：

在做实验方面，我们组一开始就落后于其他组，根系结果也不是很好，因此我们组主动申请在运动会的第二天再次去挖植物、扫描根系，这也是我们组数据偏多的原因。同时这也是我第一次在室外采集样本，确实有一定的新鲜感，让我觉得生态学实验也有一些有趣的地方

在写实验报告方面，我分析数据运用到了之前都没用过的很多方法，比如用软件RhizoVisionExplorer来扫描根系得到数据，还有主成分分析、层状聚类分析等等。这些对我来说都是很新奇的经历，即使在一开始真的有点痛苦，但也确实提升了我对数据的处理与分析的能力。

因此，我希望以后的实验也能像现在这样有趣又能让我学到不少新东西。

\section{附录}

\subsection{相关网址}

\begin{itemize}
    \item 中国植物志：http://www.cn-flora.ac.cn/
    \item 植物智：https://www.iplant.cn/
    \item PCA网站：https://www.statskingdom.com/pca-calculator.html
    \item 层次聚类分析：https://numiqo.com/statistics-calculator/cluster
\end{itemize}

\subsection{根系分析}
经Rhizo Vision Explorer处理，得到表\ref{tab:raw_data_pca_1}和表\ref{tab:raw_data_pca_2}性状的信息：

\begin{table}[htbp]
\centering
\caption{原始数据}
\label{tab:raw_data_pca_1}
\begin{tabular}{lcccc}
\toprule
文件名 & \makecell{根尖数量} & \makecell{分支点数量} & \makecell{总根长(mm)} & \makecell{每毫米分支频率} \\
\midrule
百日菊-1 & 1305 & 1799 & 2769.26 & 0.650 \\
百日菊-2 & 3857 & 5148 & 5711.89 & 0.901 \\
百日菊-3 & 4061 & 5514 & 6095.72 & 0.905 \\
合欢-1 & 1087 & 1448 & 1575.08 & 0.919 \\
合欢-2 & 858 & 1453 & 1886.93 & 0.770 \\
合欢-3 & 1061 & 1301 & 1147.50 & 1.134 \\
苎麻-1 & 4440 & 6489 & 4617.33 & 1.405 \\
苎麻-2 & 1973 & 2785 & 2657.10 & 1.048 \\
苎麻-3 & 540 & 768 & 781.61 & 0.983 \\
麦冬-1 & 4065 & 4407 & 5581.64 & 0.790 \\
麦冬-2 & 4559 & 5137 & 5793.35 & 0.887 \\
麦冬-3 & 1054 & 1148 & 1631.68 & 0.704 \\
香附子-1 & 2988 & 4906 & 4033.70 & 1.216 \\
香附子-2 & 2259 & 3601 & 3033.42 & 1.187 \\
香附子-3 & 1573 & 2787 & 2174.13 & 1.282 \\
香堇菜-1 & 1268 & 2670 & 2363.43 & 1.130 \\
香堇菜-2 & 807 & 1239 & 1257.29 & 0.985 \\
香堇菜-3 & 792 & 1930 & 1374.02 & 1.405 \\
\bottomrule
\end{tabular}
\end{table}

\begin{table}[h]
\centering
\caption{原始数据（PCA分析所用性状，表二）}
\label{tab:raw_data_pca_2}
\begin{tabular}{lccccc}
\toprule
文件名 & \makecell{网络面积(mm²)} & \makecell{平均直径(mm)} & \makecell{周长(mm)} & \makecell{体积(mm³)} & \makecell{表面积(mm²)} \\
\midrule
百日菊-1 & 1902.40 & 1.497 & 2824.23 & 10369.09 & 13576.38 \\
百日菊-2 & 3186.53 & 1.759 & 4490.54 & 34400.64 & 33518.42 \\
百日菊-3 & 3007.49 & 1.849 & 4163.03 & 41453.24 & 38662.38 \\
合欢-1 & 803.97 & 0.911 & 2032.68 & 2984.37 & 4787.46 \\
合欢-2 & 284.19 & 0.226 & 3284.06 & 205.97 & 1364.58 \\
合欢-3 & 336.59 & 0.426 & 1670.95 & 282.15 & 1571.39 \\
苎麻-1 & 1046.64 & 0.382 & 6466.17 & 1370.13 & 5712.92 \\
苎麻-2 & 725.91 & 0.469 & 3780.57 & 1182.38 & 4038.07 \\
苎麻-3 & 335.10 & 0.818 & 912.74 & 934.58 & 2104.01 \\
麦冬-1 & 3472.64 & 1.318 & 4661.58 & 14120.79 & 23837.65 \\
麦冬-2 & 2655.10 & 1.282 & 3727.41 & 13714.60 & 24117.43 \\
麦冬-3 & 1027.33 & 1.591 & 1144.77 & 5749.04 & 8384.30 \\
香附子-1 & 1381.87 & 0.946 & 4433.96 & 11842.96 & 13106.48 \\
香附子-2 & 1125.49 & 0.923 & 3208.84 & 7145.24 & 9508.45 \\
香附子-3 & 702.09 & 0.863 & 2183.56 & 3860.93 & 6295.49 \\
香堇菜-1 & 873.20 & 0.764 & 2930.70 & 4175.55 & 5920.47 \\
香堇菜-2 & 515.03 & 0.726 & 1611.58 & 1341.98 & 2946.04 \\
香堇菜-3 & 366.79 & 0.548 & 1630.59 & 981.85 & 2435.36 \\
\bottomrule
\end{tabular}
\end{table}

\newpage 

\subsection{主成分分析}

以下原始数据来自于PCA的网页，其中表头与网页一致，绘制成了紫色。

\begin{figure}[htbp]
    \centering
    \includegraphics[width=0.65\linewidth]{PCA图/ScreePlot.png}
    \caption{Scree Plot}
    \label{fig:ScreePlot}
\end{figure}

\subsubsection{特征值}

\begin{tabular}{l *{9}{c}}
\hline
\rowcolor{blue!20}
参数 & PC$_1$ & PC$_2$ & PC$_3$ & PC$_4$ & PC$_5$ & PC$_6$ & PC$_7$ & PC$_8$ & PC$_9$ \\
\hline
特征值 & 6.1677 & 1.9847 & 0.53111 & 0.1615 & 0.09175 & 0.05343 & 0.008811 & 0.0009022 & 0.00004939 \\
方差百分比 (\%) & 68.5305 & 22.0524 & 5.9007 & 1.7947 & 1.0194 & 0.5937 & 0.0979 & 0.01002 & 0.0005487 \\
累计百分比 (\%) & 68.5305 & 90.583 & 96.4837 & 98.2784 & 99.2978 & 99.8915 & 99.9894 & 99.9995 & 100 \\
\hline
\end{tabular}

\subsubsection{得分}

\begin{tabular}{l *{9}{c}}
\hline
\rowcolor{blue!20}
文件名 & PC$_1$ & PC$_2$ & PC$_3$ & PC$_4$ & PC$_5$ & PC$_6$ & PC$_7$ & PC$_8$ & PC$_9$ \\
\hline
百日菊-1 & 0.2483 & -2.0523 & -0.6965 & 0.04933 & 0.4378 & 0.3246 & -0.04917 & -0.02228 & 0.01415 \\
百日菊-2 & 4.4245 & -0.8204 & 0.5142 & -0.4013 & 0.03856 & 0.1623 & 0.09191 & 0.01549 & -0.008513 \\
百日菊-3 & 4.9377 & -1.0309 & 0.9769 & -0.6978 & -0.06016 & -0.2387 & 0.005761 & -0.02077 & 0.00714 \\
合欢-1 & -1.6573 & -0.7999 & -0.1713 & 0.0364 & 0.07335 & 0.001436 & 0.09386 & 0.008384 & 0.003364 \\
合欢-2 & -2.1234 & -0.0508 & -1.5144 & -0.8961 & -0.184 & -0.1401 & -0.1125 & 0.009112 & 0.0007691 \\
合欢-3 & -2.5842 & 0.1963 & 0.2513 & -0.112 & -0.4552 & -0.05479 & 0.1695 & -0.001113 & -0.00003818 \\
苎麻-1 & 1.2739 & 3.9046 & -0.5584 & 0.1418 & 0.4257 & -0.138 & 0.1237 & 0.00145 & 0.007402 \\
苎麻-2 & -1.0086 & 0.9585 & -0.6822 & -0.2574 & -0.01719 & 0.008721 & 0.02431 & -0.0779 & -0.01055 \\
苎麻-3 & -2.7237 & -0.9957 & 0.2966 & 0.09151 & -0.1107 & -0.16 & 0.0724 & 0.03368 & 0.002204 \\
麦冬-1 & 3.2932 & -0.3119 & -1.0358 & 0.5258 & -0.3981 & 0.477 & 0.0643 & 0.01045 & -0.002187 \\
麦冬-2 & 3.0609 & 0.07791 & -0.4449 & 0.7205 & -0.4985 & -0.4406 & -0.1422 & -0.004213 & 0.00141 \\
麦冬-3 & -1.1192 & -2.4184 & -0.0809 & 0.5146 & 0.5355 & -0.3032 & 0.03037 & -0.01948 & -0.01112 \\
香附子-1 & 1.1601 & 1.3927 & 0.2964 & -0.0821 & 0.3192 & -0.006875 & -0.07794 & 0.0404 & -0.007661 \\
香附子-2 & -0.09458 & 0.7283 & 0.4271 & 0.1468 & 0.1392 & 0.01648 & -0.03889 & 0.01427 & -0.002256 \\
香附子-3 & -1.25 & 0.5794 & 0.945 & 0.2738 & 0.03749 & 0.02985 & -0.06678 & -0.02293 & 0.009902 \\
香堇菜-1 & -1.0978 & 0.3979 & 0.1859 & -0.1101 & 0.1195 & 0.2392 & -0.1573 & 0.03163 & -0.005178 \\
香堇菜-2 & -2.2719 & -0.5882 & 0.02746 & -0.0285 & -0.1269 & -0.05547 & 0.03431 & 0.04022 & 0.003135 \\
香堇菜-3 & -2.4678 & 0.8329 & 1.2637 & 0.08477 & -0.2755 & 0.278 & -0.06567 & -0.03641 & -0.001972 \\
\hline
\end{tabular}

\subsubsection{计算步骤}

\textbf{1. 标准化每个维度的范围：}

\[
x_{\text{scaled}} = \frac{x - \text{mean}}{\text{sd}}
\]

标准化后的表格：

\begin{tabular}{l *{9}{c}}
\hline
\rowcolor{blue!20}
文件名 & 根尖数量 & 分支点数量 & 总根长 & 每毫米分支频率 & 网络面积 & 平均直径 & 周长 & 体积 & 表面积 \\
\hline
百日菊-1 & -0.5807 & -0.6802 & -0.1426 & -1.6244 & 0.5507 & 1.0991 & -0.1636 & 0.1453 & 0.2079 \\
百日菊-2 & 1.1909 & 1.1712 & 1.4858 & -0.5106 & 1.7636 & 1.6363 & 0.972 & 2.2045 & 1.9647 \\
百日菊-3 & 1.3325 & 1.3736 & 1.6982 & -0.496 & 1.5944 & 1.8209 & 0.7488 & 2.8088 & 2.4179 \\
合欢-1 & -0.732 & -0.8743 & -0.8035 & -0.4308 & -0.4868 & -0.1031 & -0.7031 & -0.4874 & -0.5663 \\
合欢-2 & -0.891 & -0.8715 & -0.6309 & -1.0915 & -0.9777 & -1.5079 & 0.1498 & -0.7255 & -0.8679 \\
合欢-3 & -0.7501 & -0.9556 & -1.0401 & 0.5184 & -0.9282 & -1.0967 & -0.9496 & -0.719 & -0.8496 \\
苎麻-1 & 1.5956 & 1.9126 & 0.8801 & 1.7204 & -0.2576 & -1.1872 & 2.3185 & -0.6258 & -0.4848 \\
苎麻-2 & -0.117 & -0.1351 & -0.2047 & 0.1394 & -0.5605 & -1.0086 & 0.4881 & -0.6418 & -0.6323 \\
苎麻-3 & -1.1118 & -1.2502 & -1.2426 & -0.1507 & -0.9296 & -0.2933 & -1.4664 & -0.6631 & -0.8027 \\
麦冬-1 & 1.3353 & 0.7616 & 1.4137 & -1.0051 & 2.0338 & 0.7313 & 1.0886 & 0.4668 & 1.1119 \\
麦冬-2 & 1.6782 & 1.1652 & 1.5309 & -0.5751 & 1.2616 & 0.6575 & 0.4519 & 0.432 & 1.1366 \\
麦冬-3 & -0.7549 & -1.0401 & -0.7721 & -1.3856 & -0.2758 & 1.293 & -1.3082 & -0.2505 & -0.2495 \\
香附子-1 & 0.5876 & 1.0374 & 0.5571 & 0.8835 & 0.05905 & -0.02982 & 0.9334 & 0.2716 & 0.1665 \\
香附子-2 & 0.08157 & 0.316 & 0.003579 & 0.7545 & -0.1831 & -0.07732 & 0.09848 & -0.1309 & -0.1504 \\
香附子-3 & -0.3947 & -0.134 & -0.4719 & 1.174 & -0.583 & -0.2018 & -0.6003 & -0.4123 & -0.4335 \\
香堇菜-1 & -0.6064 & -0.1987 & -0.3672 & 0.5004 & -0.4214 & -0.4038 & -0.09107 & -0.3854 & -0.4665 \\
香堇菜-2 & -0.9264 & -0.9898 & -0.9793 & -0.1381 & -0.7597 & -0.4819 & -0.9901 & -0.6282 & -0.7285 \\
香堇菜-3 & -0.9368 & -0.6078 & -0.9147 & 1.7172 & -0.8997 & -0.8469 & -0.9771 & -0.659 & -0.7735 \\
\hline
\end{tabular}

\textbf{2. 计算协方差矩阵：}

\[
\text{Cov}_{x_1,x_2} = \frac{\sum (x_{1,i} - \bar{x}_1)(x_{2,i} - \bar{x}_2)}{n - 1}
\]

协方差矩阵：

\begin{tabular}{p{1.85cm} *{9}{c}}
\hline
\rowcolor{blue!20}
Group$\backslash$Group & 根尖数量 & 分支点数量 & 总根长 & 每毫米分支频率 & 网络面积 & 平均直径 & 周长 & 体积 & 表面积 \\
\hline
根尖数量 & 1 & 0.9567 & 0.9645 & 0.03186 & 0.8006 & 0.4219 & 0.8528 & 0.6425 & 0.7634 \\
分支点数量 & 0.9567 & 1 & 0.9167 & 0.2421 & 0.6801 & 0.3163 & 0.8976 & 0.607 & 0.6826 \\
总根长 & 0.9645 & 0.9167 & 1 & -0.1242 & 0.9009 & 0.5637 & 0.8177 & 0.7792 & 0.8809 \\
每毫米分支频率 & 0.03186 & 0.2421 & -0.1242 & 1 & -0.4106 & -0.5504 & 0.1405 & -0.3025 & -0.3502 \\
网络面积 & 0.8006 & 0.6801 & 0.9009 & -0.4106 & 1 & 0.7923 & 0.5641 & 0.8396 & 0.9455 \\
平均直径 & 0.4219 & 0.3163 & 0.5637 & -0.5504 & 0.7923 & 1 & 0.0986 & 0.8119 & 0.8306 \\
周长 & 0.8528 & 0.8976 & 0.8177 & 0.1405 & 0.5641 & 0.0986 & 1 & 0.4216 & 0.4986 \\
体积 & 0.6425 & 0.607 & 0.7792 & -0.3025 & 0.8396 & 0.8119 & 0.4216 & 1 & 0.9632 \\
表面积 & 0.7634 & 0.6826 & 0.8809 & -0.3502 & 0.9455 & 0.8306 & 0.4986 & 0.9632 & 1 \\
\hline
\end{tabular}

\textbf{特征向量：}

\begin{tabular}{*{9}{c}}
\hline
\rowcolor{blue!20}
Vector1 & Vector2 & Vector3 & Vector4 & Vector5 & Vector6 & Vector7 & Vector8 & Vector9 \\
\hline
0.3677 & 0.2449 & -0.1208 & 0.3519 & -0.2036 & -0.4777 & 0.6064 & -0.1646 & -0.0438 \\
0.3442 & 0.3593 & 0.06217 & 0.07771 & 0.224 & -0.2782 & -0.3308 & 0.6903 & 0.1705 \\
0.3941 & 0.1182 & -0.1085 & 0.07041 & -0.1971 & -0.06875 & -0.6112 & -0.3908 & -0.494 \\
-0.09027 & 0.5796 & 0.7203 & 0.1339 & -0.02617 & 0.2866 & 0.05246 & -0.1818 & -0.02294 \\
0.3811 & -0.1574 & -0.09206 & 0.3018 & -0.345 & 0.6816 & 0.1476 & 0.3439 & -0.08233 \\
0.2846 & -0.4429 & 0.2878 & 0.4034 & 0.6668 & 0.01717 & 0.01676 & -0.1756 & -0.03494 \\
0.2931 & 0.407 & -0.4223 & -0.3466 & 0.4673 & 0.3637 & 0.1624 & -0.2392 & 0.1258 \\
0.3527 & -0.2048 & 0.3795 & -0.6668 & -0.02844 & -0.09562 & 0.2631 & 0.1608 & -0.3741 \\
0.3826 & -0.1804 & 0.1931 & -0.1695 & -0.2933 & -0.03429 & -0.1697 & -0.2801 & 0.7487 \\
\hline
\end{tabular}

\subsection{层次聚类分析}

该网页只给出了肘部法则图，如下：

\begin{figure}[htbp]
    \centering
    \includegraphics[width=0.7\linewidth]{树状/Elbow Method.png}
    \caption{肘部法则}
    \label{fig:placeholder}
\end{figure}

\end{document}